\documentclass[12pt]{article}
\usepackage{amsmath,amssymb,amsfonts,amsthm}
\usepackage[margin=1in]{geometry}

\begin{document}

\begin{center}
{\Large\bfseries Chapter 9, Problem 10}\\[6pt]
Byron \& Fuller, \textit{Mathematics of Classical and Quantum Physics}
\end{center}

\bigskip

\textbf{Problem.} Consider the quantity
\[
\Phi = \langle f, (A_0 - \lambda_1^{(0)})f\rangle + 2\langle f, (A_1 - \lambda_1^{(1)})x_1^{(0)}\rangle,
\]
where $\lambda_1^{(0)}$ is the largest eigenvalue of $A_0$ with eigenvector $x_1^{(0)}$, $A_1$ is a small perturbation, $\lambda_1^{(1)} = \langle x_1^{(0)}, A_1 x_1^{(0)}\rangle$, and the spectrum of $A_0$ is nondegenerate. Show that $\Phi \le \lambda_1^{(2)}$ for all $f \in H$, with equality iff $f = f^{(1)}$ as given by Eq.~(9.40).

\bigskip
\textbf{Solution.}

We work in a real Hilbert space as specified. Expand $f$ in eigenvectors of $A_0$:
\[
f = \sum_n c_n\,x_n^{(0)}, \quad c_n = \langle x_n^{(0)}, f\rangle.
\]

\textbf{First term:} Since $A_0 x_n^{(0)} = \lambda_n^{(0)} x_n^{(0)}$:
\[
\langle f, (A_0 - \lambda_1^{(0)})f\rangle = \sum_n c_n^2(\lambda_n^{(0)} - \lambda_1^{(0)}).
\]
Since $\lambda_1^{(0)}$ is the \emph{largest} eigenvalue, $\lambda_n^{(0)} - \lambda_1^{(0)} \le 0$ for all $n$, and $= 0$ only for $n = 1$.

\textbf{Second term:} Define $V_{n1} = \langle x_n^{(0)}, A_1 x_1^{(0)}\rangle$. Note $V_{11} = \lambda_1^{(1)}$. Then:
\[
\langle f, (A_1 - \lambda_1^{(1)})x_1^{(0)}\rangle = \sum_n c_n(V_{n1} - \lambda_1^{(1)}\delta_{n1}) = \sum_{n \neq 1} c_n V_{n1}.
\]

Combining:
\[
\Phi = \sum_n c_n^2(\lambda_n^{(0)} - \lambda_1^{(0)}) + 2\sum_{n \neq 1} c_n V_{n1}.
\]

The $n = 1$ term in the first sum vanishes, so:
\[
\Phi = \sum_{n \neq 1} \left[c_n^2(\lambda_n^{(0)} - \lambda_1^{(0)}) + 2c_n V_{n1}\right].
\]

For each $n \neq 1$, complete the square. Let $\Delta_n = \lambda_n^{(0)} - \lambda_1^{(0)} < 0$:
\[
\Delta_n c_n^2 + 2c_n V_{n1} = \Delta_n\!\left(c_n + \frac{V_{n1}}{\Delta_n}\right)^{\!2} - \frac{V_{n1}^2}{\Delta_n}.
\]

Since $\Delta_n < 0$, the squared term $\Delta_n(c_n + V_{n1}/\Delta_n)^2 \le 0$. Therefore:
\[
\Phi = \sum_{n \neq 1}\left[\Delta_n\!\left(c_n + \frac{V_{n1}}{\Delta_n}\right)^{\!2} - \frac{V_{n1}^2}{\Delta_n}\right] \le -\sum_{n \neq 1} \frac{V_{n1}^2}{\Delta_n} = \sum_{n \neq 1} \frac{V_{n1}^2}{\lambda_1^{(0)} - \lambda_n^{(0)}}.
\]

Now, from Eq.~(9.42), the second-order perturbation correction is:
\[
\lambda_1^{(2)} = \sum_{n \neq 1} \frac{|\langle x_n^{(0)}, A_1 x_1^{(0)}\rangle|^2}{\lambda_1^{(0)} - \lambda_n^{(0)}} = \sum_{n \neq 1} \frac{V_{n1}^2}{\lambda_1^{(0)} - \lambda_n^{(0)}}.
\]

Therefore:
\[
\boxed{\Phi \le \lambda_1^{(2)}}.
\]

\textbf{Equality} holds when each squared term vanishes, i.e., $c_n = -V_{n1}/\Delta_n = V_{n1}/(\lambda_1^{(0)} - \lambda_n^{(0)})$ for all $n \neq 1$ (and $c_1$ is arbitrary since the $n=1$ term cancels). This means:
\[
f = c_1 x_1^{(0)} + \sum_{n \neq 1} \frac{V_{n1}}{\lambda_1^{(0)} - \lambda_n^{(0)}}\,x_n^{(0)}.
\]

The sum (ignoring the $c_1 x_1^{(0)}$ component, which doesn't affect $\Phi$) is precisely $f^{(1)}$ from Eq.~(9.40):
\[
f^{(1)} = \sum_{n \neq 1} \frac{\langle x_n^{(0)}, A_1 x_1^{(0)}\rangle}{\lambda_1^{(0)} - \lambda_n^{(0)}}\,x_n^{(0)}.
\]

\textbf{Extension to complex Hilbert space:} Replace $\Phi$ by
\[
\Phi = \langle f, (A_0 - \lambda_1^{(0)})f\rangle + 2\,\text{Re}\,\langle f, (A_1 - \lambda_1^{(1)})x_1^{(0)}\rangle,
\]
since all quantities must be real for the variational principle to apply. \hfill $\blacksquare$

\end{document}
